\section{Educational Context and Learning Objectives}

In today's tech-driven world, there's a growing need to develop skills for navigating the digital era. 
Computational thinking (CT) has become vital, fostering analytical and problem-solving skills crucial for coding proficiency \cite{tsai_developing_2019}. 
Integrating programming skills into education can enhance CT abilities, that are essential for future readiness, particularly among primary school students \cite{popat_learning_2019}.

\subsection{Educational programming tools for children}

Typically, children begin learning programming at the stage of middle school. 
One of the tools they are taught with, is Scratch, a popular simplified block-based language. 
Scratch's aesthetic aspects are the most significant component for developing computational thinking \cite{oscco_solorzano_analysis_2025}. 
Research indicates that game-based programming environments like Scratch serve as crucial factors for improving education quality by changing student attitudes toward the frustrations of learning programming logic and fostering skill development through experimentation in motivating contexts \cite{kuz_computational_2023}. 
Other effective tools for teaching programming are Kodable, CodeMonkey, BlocklyGames. 
Their main effectiveness in developing computational thinking comes from their gamified design, which engages children by providing immediate feedback and having a memorable interface \cite{hlazova_integration_2024}.

\subsection{Gamification of learning programming}

The integration of game elements into the learning process can contribute to the development of algorithmic thinking, logic and creative abilities of students from an early age \cite{medvedieva_use_2021}. 
Game methods of teaching programming differ from others in that the game is a natural learning environment that attracts students, increases their motivation and stimulates active research, allows students to apply theoretical knowledge in practice, receiving immediate feedback and observing the results of their work, promotes the development of not only technical programming skills, but also such important qualities as creativity, critical thinking and problem-solving skills, and is an interactive process that allows students to interact with educational material in a dynamic and exciting environment, creates a center for cooperation and exchange of experience between students, contributing to the development of team skills \cite{moshkova_gamification_2024}. 
The perceived enjoyment and social influence yielded by gamified programming instruction have a high impact on student's motivation to learn programming \cite{faculty_of_education_in_jagodina_university_of_kragujevac_serbia_exploring_2025}.

\subsection{Educational DSL Precedents and Benefits}

Logo represents one of the earliest educational DSLs, designed in 1967 to teach programming concepts through turtle graphics and LISP-like constructs \cite{noauthor_domain-specific_nodate}. 
Game Maker Language similarly combines simplified commands derived from GPLs (General Purpose Languages) to enable users to create games without advanced programming knowledge \cite{noauthor_domain-specific_nodate}. 
Well-designed DSLs improve programmer productivity, addressing key challenges in software development \cite{fowler_domain-specific_2011}.
